\question
После долгих размышлений у печки и научного спора с Дядей Фёдором, рациональный кот Матроскин решил подойти к вопросу бутербродоедения с истинно научной строгостью. Он скрупулёзно записывал сведения о бутебродах и их вкусовых качествах. Результатом наблюдений стало то, что сэндвич считается идеальным если выполняется одно из условий: (колбаса сверху И хлеб свежий И (масло ИЛИ икра)) ИЛИ (колбаса снизу И (свежий хлеб XOR масло) И НЕ хлеб черствый). В остальных случаях бутерброд получался "так себе". 

Матроскин недавно увидел статью про карты Вейча, как способ минимизации функций, но, к сожалению, ничего не понял, поэтому решил попросить вашей помощи.

Превратите эти гастрономические наблюдения в логическую формулу, да не простую, а компактную (найти минимальную ДНФ), чтобы Матроскин мог быстро определять оптимальный способ бутербродоукладывания.

---------------

Автор -- Никита Глейм